% Shared visual style with the Purshow Notes collection.
\documentclass[11pt,fontset=none]{ctexart}
\usepackage[a3paper,landscape,left=15mm,right=15mm,top=19mm,bottom=17mm,headheight=15pt,headsep=9pt,footskip=22pt]{geometry}
\usepackage{fontspec,amsmath,amssymb,array,longtable,booktabs,fancyhdr,xcolor,xurl,hyperref,graphicx,enumitem,needspace}
\setmainfont[Path=/usr/local/texlive/2026/texmf-dist/fonts/opentype/public/tex-gyre/,UprightFont=texgyrepagella-regular.otf,BoldFont=texgyrepagella-bold.otf,ItalicFont=texgyrepagella-italic.otf,BoldItalicFont=texgyrepagella-bolditalic.otf]{texgyrepagella-regular.otf}
\setsansfont[Path=/usr/local/texlive/2026/texmf-dist/fonts/opentype/public/tex-gyre/,UprightFont=texgyreheros-regular.otf,BoldFont=texgyreheros-bold.otf,ItalicFont=texgyreheros-italic.otf,BoldItalicFont=texgyreheros-bolditalic.otf]{texgyreheros-regular.otf}
\setmonofont[Path=/usr/local/texlive/2026/texmf-dist/fonts/opentype/SIL/jetbrainsmono-otf/,UprightFont=JetBrainsMono-Regular.otf,BoldFont=JetBrainsMono-Bold.otf,Scale=0.82]{JetBrainsMono-Regular.otf}
\setCJKmainfont[Path=/System/Library/Fonts/Supplemental/,FontIndex=6,BoldFont=Songti.ttc,BoldFeatures={FontIndex=1}]{Songti.ttc}
\setCJKsansfont[Path=/System/Library/Fonts/,FontIndex=1,BoldFont={STHeiti Medium.ttc},BoldFeatures={FontIndex=1}]{STHeiti Light.ttc}
\setCJKmonofont[Path=/System/Library/Fonts/,FontIndex=1]{STHeiti Light.ttc}
\definecolor{notesblue}{HTML}{174A7E}
\definecolor{notesgray}{HTML}{5D6773}
\definecolor{noteslink}{HTML}{1769AA}
\hypersetup{colorlinks=true,urlcolor=noteslink,linkcolor=noteslink,pdftitle={视觉编码器架构与多模态训练对比},pdfauthor={Purshow},pdfsubject={视觉模块、时空压缩、位置编码、动态分辨率、多模态训练 Recipe、输入 infra}}
\setlength{\parindent}{0pt}
\setlength{\parskip}{3pt}
\pagestyle{fancy}
\fancyhf{}
\fancyhead[L]{\small\sffamily\color{notesgray}Vision Encoder / Architecture Notes}
\fancyhead[R]{\small\sffamily\color{notesgray}Purshow's Notes}
\fancyfoot[R]{\small\thepage}
\renewcommand{\headrulewidth}{0.3pt}
\setcounter{secnumdepth}{0}
\raggedbottom
\newcolumntype{P}[1]{>{\raggedright\arraybackslash}p{#1}}
